\subsection{Potential Change as Tangents}

Differential Calculus tracks how functions adapt to changes in
their inputs. The tangent space at a point is the vector space of
all potential changes applicable at that point. The forward and
reverse derivatives transport changes forward and back. We track
data provenance using a simplification of this idea where tangents
are taken from the Boolean semiring $\{0,1\}$ where $0$ represents
``no change'' and $1$ represents ``possible change''. Potential
change to an $n$-tuple of data is represented as an $n$-vector of
$\{0,1\}$ values. These vectors form a semimodule over the semiring
structure on $\{0,1\}$ with maximum (join) as addition and minimum
(meet) as multiplication. Thus ``no change'' is represented as the
$0$ vector $(0 \dots 0)$.

With this definition of ``tangent'', the derivative of a
function
$f : A_1 \times \cdots \times A_m \to B_1 \times \cdots \times B_n$
ought to map tangent vectors in $\{0,1\}^m$ to tangent vectors in
$\{0,1\}^n$, describing how potential changes in the input are
reflected in the output. Differentiable functions on manifolds have
an intrinsic notion of derivative, but for plain functions $f$ we
require that they are paired with a tangent map that describes the
intensional information about how they depend on their
arguments. For $f$, the tangent map $\partial f$ has type
$A_1 \times \cdots \times A_n \to \{0,1\}^{m\times n}$, mapping
points to matrices representing linear maps: the ``Jacobian'' for
this function. We are not free to choose an arbitrary linear map for
$\partial f(a_1, \dots, a_m)$, but must choose one that satisfies
the following \emph{dependency correctness} property:

\begin{definition}[Dependency Correctness] A tangent map
  $\partial f : \overrightarrow{A_i} \to \{0,1\}^{m \times n}$ is
  \emph{dependency correct} for a function
  $f : \overrightarrow{A_i} \to \overrightarrow{B_j}$ if for all
  $a, a' \in \overrightarrow{A_i}$ and $0 \leq j < m$, if $a$ and
  $a'$ are equal on all components $i$ such that
  $\partial f(a)_{i,j} = 1$, then $f(a)_j = f(a')_j$.
\end{definition}

Reading the definition backwards, this says that $f$ is
constant at component $j$ on inputs that differ only at components
$i$ where the tangent map matrix says $0$. Thus we read the entry at
$(i,j)$ of the tangent matrix as capturing the dependency of output
$j$ on input $i$. We illustrate this with an example:

\begin{example}
  Let $\pi_1 : A_1 \times A_2 \to A_1$ be first projection. At all
  points $(a_1, a_2)$, the tangent matrix is $(1\;0)$ indicating
  that the output depends on the first argument but not the
  second. This is justified by observing how it acts on tangent
  vectors. Representing no change in the first argument and
  potential change in the second as
  $\left(\begin{array}{c} 0 \\ 1 \end{array}\right)$, applying the
  tangent map gives $0$ indicating no change in the output. Swapping
  the potential change to
  $\left(\begin{array}{c} 1 \\ 0 \end{array}\right)$ yields the
  output tangent $1$ indicating that changing the first argument
  changes the output. By linearity, all other potential changes are
  derivable from these. In particular, for all tangent maps, no
  change in the input ($\overrightarrow{0}$) is always mapped to no
  change in the output.
\end{example}

Note that dependency correct tangent maps for a function are
not unique. We can always overapproximate the dependency of a
function on its input. In the case of a binary function like first
projection, $\partial\pi_1(a_1,a_2) = (1\;1)$ is also permissble.

\begin{example}
  Multiplication $(*) : \mathbb{N} \times \mathbb{N} \to \mathbb{N}$
  is an example of a function that has tangent maps that depend on
  the input. Dependency correct tangent maps can be defined as:
  \begin{displaymath}
    \begin{array}{lcl}
      \partial(*)(0,0)&=&(1\;1) \\
      \partial(*)(0,n)&=&(1\;0) \\
      \partial(*)(m,0)&=&(0\;1) \\
      \partial(*)(m,n)&=&(1\;1)
    \end{array}
  \end{displaymath}
  These maps indicate that the dependency of multiplication on an
  input depends on whether or not other other input is $0$. If the
  first argument is zero, then changing the second will have no
  effect and vice versa. If both arguments are $0$, then a change in
  both will change the output. However, since tangent maps are
  componentwise we have no way of saying change in both, so we
  overapproximate by saying that a change in either \emph{may}
  change the output. Again, this choice is not unique. We could also
  use $(1\;0)$ or $(0\;1)$ and still be dependency correct,
  indicating a choice to be left- or right-biased.
\end{example}

Functions with dependency correct tangent maps have an identity
and compose according to the chain rule, so they form a category.

\newcommand{\DepCat}{\mathbf{DepCat}}

\begin{proposition}
  The identity tangent map is dependency correct for the identity
  function. Composable functions $f, g$ with dependency correct
  tangent maps $\partial f$ and $\partial g$ are closed under
  composition, with the tangent map being
  $\partial(f \circ g)(x) = \partial f (g(x)) \circ \partial g(x)$.

  The category $\DepCat$ of functions with dependency information
  has objects that are tuples of sets $\overrightarrow{A_i}$ and
  morphisms that are functions with dependency correct tangent
  maps. This category has finite products.
\end{proposition}

FIXME: finite products is interesting, because it determines the
dependency information???

The fact that dependency correct functions form a category allows us
to model a simple programming language with finite product
types. Primitive operations can be interpreted a functions that
overapproximate their dependency information:

\begin{proposition}
  There is a functor $\Set \to \DepCat$ that preserves finite
  products. FIXME: check!
\end{proposition}

Morphisms in $\DepCat$ constructed from finite product structure and
embedded functions will only have trivial tangent maps arising from
the use of projections. A generator of non-trivial dependency
information is the ability to use parts of the input depending on
the input itself, via a condition (if-then-else) map:

\begin{proposition}
  Let $\mathbb{B}$ be the object of $\DepCat$ consisting of the set
  of booleans $\{\mathsf{true}, \mathsf{false}\}$. There are
  dependency correct maps
  $\mathrm{true}, \mathrm{false} : 1 \to \mathbb{B}$ and a natural
  family of maps
  $\mathrm{cond}_X : \mathbb{B} \times (X \times X) \to X$ such that $\mathrm{cond}_X \circ \langle \mathrm{true} \circ !, \mathrm{id} \rangle = \pi_1$ and $\mathrm{cond}_X \circ \langle \mathrm{false} \circ !, \mathrm{id} \rangle = \pi_2$. FIXME: and extensionality.
\end{proposition}

Thus $\DepCat$ is suitable for modelling a small language with
booleans, conditionals, and a collection of primitive types with
operations. This setting does not cover our example from the
introduction, which used lists and higher-order functions. In
Section \ref{sec:models-of-total-gps} we embed $\DepCat$ into
categories that account for sum types, some algebraic types and
higher-order functions.
